\subsection{Background and history}\label{sub:history}

Integer Cantor sets, also called ``restricted digit sets'', are a natural object of study.  The earliest result 
known to us is from 
1965 in Fine's paper \cite{MR184895}, which studied the distribution of the sum of digits function modulo $p$, a question closely related to Proposition \ref{p;sumDigits} below. 
Weyl's polynomial equidistribution theorem has been extended to Cantor sets by Coquet  \cites{MR506123,MR578819}. We will present an extension of these results below.  This was followed by work of Erd\H{o}s, Mauduit and S\'{a}rk\"{o}zy \cites{ERDOS199899,MR1692287} who considered the modular joint  distribution of Cantor sets of integers, with their   `sum of digits' function.  They established uniform distribution modulo $n$, with error rates on convergence.  
Their work is extended by 
Konyagin \cite{MR1832701}
and more recently 
Saavedra-Araya in \cite{SAAVEDRA-ARAYA_2026}, who used Markov chain techniques to generalize  such sets.  

The definition of `classical' Cantor sets has many possible extensions, with definitions in terms of systems of linear equations \cites{MR2678855,MR2290789}, 
automata, as for instance the Thue--Morse sequence \cite{MR1830572}, Markov Chains \cite{SAAVEDRA-ARAYA_2026},
and more general dynamical variants 
as studied by 
Glasscock, Moreira and Richter 
\cite{MR4734865}. 
 


The role of digit functions and restricted digits along the prime integers was raised by Gelfond \cite{MR220693}, who asked for the distribution of the sum of digits function, for a given base $b$, mod $m$. This was answered 
by Mauduit and Rivat \cite{MR2680394}.
This result attracted a great deal of attention to these and related problems. 


Particularly relevant is Maynard's seminal work \cite{MaynardInvent2019}, which proved the existence of infinitely many prime numbers in certain kinds of Cantor sets; the work of Leng and Sawhney \cite{leng2024vinogradovstheoremprimesrestricted}, which established a Goldbach-type result, extending Vinogradov's theorem to this setting; and the work of Nath \cite{jlms.12837}, who proved Bombieri--Vinogradov type theorems for primes with a missing digit. 
In a different direction Kra and Shalom \cite{kra2025ergodicaverageslargeintersection}  
have studied multiple recurrence for `rational spectra' IP sets.  Here IP can stand for   `infinite parallelepipeds'.  In our setting, any Cantor set with base $b$ and two allowed digits, one of which is zero, forms an IP set, with rational spectra.   

